\documentstyle[%


righttag%


,11pt%


,amssymb%


,russian%               remove in Eng.


,izvru%                 Set izven% in Eng.


% ,izvru-ks             for brief communications


]{amsart}





% repl. {,} for {.}





\newcommand{\im}{\operatorname{Im}}


\newcommand{\re}{\operatorname{Re}}


\newcommand{\const}{\operatorname{const}}


\newcommand{\tts}[1]{}





\theoremstyle{plain}


\theorembodyfont{\it}


\newtheorem{thmnonum}{\hskip\parindent Теорема}


\renewcommand{\thethmnonum}{}








%\hoffset=-15mm%                  For Eng.


\setcounter{page}{14}


\god{2005}


\nomer{10~(521)} %                        Remove in Eng.


%\nomer{Vol.\,49, No.\,10}%               Set in Eng.


%\str{\pageref{begin}--\pageref{end}}%  Set in Eng.


\udk{517.518}





\author{\it Ф.Н.\,Гарифьянов, Б.\,Туре}


% NB:   ^^ remove in Eng.


\title{Проблема обращения особого 


интеграла\\ на круговом секторе}





\date{}


%\pagestyle{myheadings}%         Set in Eng.


\pagestyle{plain} %              Remove in Eng.


\begin{document}


%\thispagestyle{plain}%          Set in Eng.


\maketitle


\label{begin}














Пусть $D$ --- круговой сектор, ограниченный


отрезками $l_1$ и $l_3$ прямых


$\re z=\mp \im z$ соответственно и дугой $l_2$ окружности


$|z|=1$, $|\arg z|< \pi/4$.


Если из замыкания $\overline{D}$ удалить часть положительно


ориентированной границы $\Gamma=\partial D$,  то получим


фундаментальное множество частного случая конечной группы диэдра


([1], гл.\,7, \S\,6). Порoждающие преобразования


$\sigma_{1}(z)=iz$, $\sigma_{2}(z)=z^{-1}$ индуцируют гомеоморфизм


$\alpha (t)=\Gamma \to \Gamma$,


$$


\alpha (t)=\overline{t}\ \Longleftrightarrow\ \alpha (t)=\{


\sigma_j(t),\ t\in l_j\},


\ \ j=1,2,3; \quad\sigma_3=\sigma_1^{-1} ;\quad\alpha


(\alpha (t))\equiv t ,


$$


изменяющий ориентацию границы. Производная $\alpha'(t)$


разрывна в вершинах $t_1=0$, $t_{2,3}=\sqrt{2}/2\mp i\sqrt{2}/2$,


а сам сдвиг $\alpha(t)\in C(\Gamma)$.





Целью данной работы является исследовaние проблемы обращения особого


интеграла


\begin{equation}


(A\varphi)(t)\equiv\frac{1}{\pi


i}\int_{\Gamma}\varphi(\tau)E(t,\tau)d\tau =\psi(t),\quad


t\in\Gamma,


\end{equation}


где функции $\varphi$ и $\psi$


удовлетворяют условию Г\"eльдера на $\Gamma$. Класс таких функций


обозначим через $Q$. Ядро интеграла (1) определяется формулой


\begin{equation}


E(z,\tau)=(\tau-iz)^{-1}+(\tau+iz)^{-1}+(\tau-z^{-1})^{-1}.


\end{equation}


Здесь интеграл (1) понимается в смысле главного значения по Коши.


Заметим, что для соответствующей функции $(A\varphi)(z)$


справедлив аналог формулы Ю.В.\,Сохоцкого--Й.\,Племеля


\begin{equation}


A^{+}=-W+A,


\end{equation}


где инволютивный оператор $W$ определяется формулой


$W:\varphi(t)\to\varphi[\alpha(t)]$. Очевидно, что для


граничного значения аналитической в $D$ функции $a(z)$ имеем


\begin{equation}


Aa^{+}=Wa^{+},


\end{equation}


т.\,к. преобразования группы, отличные от


тождественного, переводят точку из области $D$ в ее внешность.


Займемся регуляризацией уравнения (1). Для этого заметим, что (1)


$\Longrightarrow A^{2}\varphi=A\psi$ и в силу формул (3) и (4)


справедливо равенство $A^{2}=A[-W+A+W]=-J+WA+AW$. Ядро оператора


$AW+WA$ имеет вид


\begin{equation}


K(t,\tau)=E(t,\tau)-\alpha'(\tau)E(\alpha(t),\alpha(\tau)).


\end{equation}


Оно ограничено, что непосредственно проверяется перебором


различных вариантов взаимного расположения точек $\tau$ и $t$ на


сторонах $\Gamma$. Более того, частные производные $\partial


K/\partial \tau$ и $\partial K/\partial t$ непрерывны на замыкании


каждой из сторон, а вершины являются для них точками разрыва


первого рода. Поэтому операторы $A$ и $W$


антикоммутируют (с точностью до компактного оператора)


в целом ряде функциональных пространств. Другими словами,


оператор $A^2=-I+F$, где~$F$ вполне непрерывен.


Некоторые вопросы теории таких операторов изучались ранее


в~[2].








Очевидно, $A: Q\longrightarrow Q$. Введем оператор


\begin{equation}


T_{\lambda}=\lambda I+AW+WA,\quad\lambda\ne 0.


\end{equation}


Проблема обращения (1) свелась к исследованию интегрального


уравнения


\begin{equation}


T_{-1} \varphi=A\psi.


\end{equation}


Фундаментальную систему решений (ф.\,с.\,р.) однородного уравнения


\begin{equation}


T_{\lambda} \varphi=0


\end{equation}


всегда можно выбрать так, чтобы каждая входящая в нее функция


удовлетворяла одному из условий


\begin{equation}


\varphi=W\varphi\quad (\varphi=-W\varphi).


\end{equation}


Такие решения будем называть, как обычно, инвариантными


(антиинвариантными). Если уравнение (8) имеет нетривиальное


решение, то оно имеет нетривиальное решение одного из видов (9)


или обоих видов, если первоначально решение не удовлетворяет ни


одному из условий (9).








\begin{lem}


Если


$\lambda\neq\pm 2$ и ф.\,с.\,р. уравнения $(8)$ содержит функцию с


одним из свойств $(9)$, то она содержит и функцию с противоположным


свойством.


\end{lem}








\begin{pf}


Пусть для определенности решение $\varphi$  инвариантно. Запишем


уравнение (8) в виде


$-\lambda\varphi=WA\varphi+A\varphi \Longleftrightarrow


-(\lambda+2)\varphi=A^{+}\varphi+WA^{+}\varphi$.


Это означает, что функция


\begin{equation}


\theta=-2^{-1}(\lambda+2)\varphi-A^{+}\varphi


\end{equation}


антиинвариантна. Умножим обе части равенства (10) на ядро (2) и


проинтегрируем по $\Gamma$, откуда


$A^{+}\theta=-2^{-1}(\lambda+2)A^{+}\varphi$. Используя формулу


(10) и исключая величину $A^{+}\varphi$, получим


$A^{+}\theta=2^{-1}(\lambda+2)[\theta+2^{-1}(\lambda+2)\varphi]$.


Применим к обеим частям последнего равенства оператор $W$, а


затем полученное соотношение вычтем из исходного. Тогда


$A^{+}\theta-WA^{+}\theta=(\lambda+2)\theta \Longleftrightarrow


(8)$. Допустим, что $\theta\equiv 0$. В силу (10) имеем


$\varphi=\varphi^{+}$, т.\,е. с учетом инвариантности


$\varphi\equiv \const$. Но постоянная принадлежит ф.\,с.\,р. уравнения


(8) только при $\lambda=-2$.


\end{pf}








\begin{lem}


Однородное уравнение Фредгольма


$(8)$ при $\lambda=-1$ не имеет нетривиальных решений.


\end{lem}








\begin{pf}


В силу леммы 1 достаточно


показать, что ф.\,с.\,р. уравнения $T_{-1}


\varphi=0$ не содержит антиинвариантных функций. Рассмотрим


это уравнение в банаховом пространстве


$C(\Gamma)$ с нормой $\|\varphi\|=\max|\varphi(t)|$, $t\in\Gamma$.


Возможны два случая.





{\bf 1${}^\circ$}. $t\in l_{1}$. Разобьем


интегральное слагаемое в (6) на три интеграла $I_{j}$ по


соответствующим сторонам $\Gamma$. При $\tau\in l_{1}$ имеем


$K(t,\tau)=2t(t^{2}\tau^{2}-1)^{-1}$, т.\,е.


$|I_{1}|<2\pi^{-1}\|\varphi\|$. При $\tau\in l_{3}$ имеем


$K(t,\tau)\equiv 0$. При $\tau\in l_{2}$ получим


$K(t,\tau)=2\tau^{-1}+(\tau+it)^{-1}+(\tau(t\tau +1))^{-1}$. За


счет антиинвариантной плотности справедливо равенство


$$


2\pi iI_{2}=\int_{l_{2}}\varphi(\tau)[K(t,\tau)-


\tau^{-2}K(t,\tau^{-1})]d\tau,


$$


в чем можно убедиться заменой переменой $\tau=\tau_{1}^{-1}$ в


интеграле от второго слагаемого в ядре. Выполняются очевидные


неравенства


$$ 


|(t\tau^{2}+\tau)^{-1}-(\tau+t)^{-1}|<0{,}5\,,\quad


|(\tau+it)^{-1}-(it\tau^{2}+\tau)^{-1}|<0{,}5\,,


$$


т.\,е. $|I_{2}|\le 0{,}25\|\varphi\|$.  Из полученных оценок следует


$\|\varphi\|=0$.





{\bf 2${}^\circ$}. $t\in l_{2}$. При


$\tau\in l_{2}$ имеем $K(t,\tau)=3\tau^{-1}$ , т.\,е.


$I_{2}=0$. Кроме того, 


$$


K(t,\tau)=\frac{1}{\tau-it}-


\frac{t}{t\tau+1}=\frac{1+it^{2}}{(\tau-it)(t\tau+1)}


\quad \text{при $\tau\in l_{1}$}


$$


и 


$$


K(t,\tau)=\frac{1}{\tau+it}-\frac{t}{t\tau+1}=


\frac{1-it^{2}}{(\tau+it)(t\tau+1)}\quad\text{при $\tau\in l_{3}$}.


$$


Поскольку


$$


|\tau-it|\,|t\tau+1|\geq 1\ \ (\tau\in l_{1})\quad


\text{и}


\quad |\tau+it|\,|t\tau+1|\geq 1\ \ (\tau\in l_{3}),


$$


то


$$


|I_{1}+I_{3}|\le \frac{1}{\pi}\max \{ |1+it^{2}|+|1-it^{2}|\}


\|\varphi\|\le\frac{2\sqrt{2}}{\pi}\|\varphi\|


\ \Longrightarrow \ \varphi\equiv 0.


$$


Тем самым доказано, что однородное уравнение не имеет


нетривиальных непрерывных решений.


\end{pf}





Сформулируем основной результат.








\begin{thmnonum}


Проблема обращения особого


интеграла $(1)$ с ядром $(2)$, понимаемого в смысле главного значения


по Коши, безусловно разрешима и имеет единственное решение


$$


\varphi=T^{-1}_{-1}(A\psi).


$$


\end{thmnonum}





Для доказательства нужно применить альтернативу Фредгольма и лемму 2.





\begin{thebibliography}{9}





\bibitem{1}


Голубев В.В. {\em Лекции по аналитической теории


дифференциальных уравнений}. -- М.--Л.: ГИТТЛ, 1950. -- 436~с.


\bibitem{2}


Агаев Г.Н. {\em К теории сингулярного уравнения в


пространствах Банаха} // Тр. ин-та физ. и мат. АН АзССР. Cер. матем.


-- 1959. -- T.\,8. -- C.\,23--27.





\end{thebibliography}





\vskip 2cm





\post{\it Казанский государственный}{\it Поступила}


\post{\it энергетический университет}{27.05.2004}


\smallskip


\post{\it Казанский государственный}{}


\post{\it университет}{}





\label{end}


\end{document}





